% !TEX root = ../main.tex
% Chapter 3 --- Data and Materials
% Self-contained section file: no preamble, no \documentclass.
% Requires the document preamble to load: booktabs, tikz, biblatex/natbib.
% TikZ libraries assumed available: positioning, arrows.meta, shapes.geometric, fit, backgrounds, calc.

\chapter{Data and Materials}
\label{ch:data}

This chapter documents every data source the project touches, the role each
plays, and the data-hygiene facts that constrain what can honestly be claimed
from them. It deliberately mirrors the template's data-acquisition chapter
(collection flow, inclusion/exclusion, expert re-check) but adapts that
structure to a \emph{multi-source} reality in which a single in-domain binary
corpus is surrounded by scaffolding, preprocessing, and unit-test fixtures that
must never be allowed to co-mingle. Two framing rules govern the entire
chapter and are stated once here so that they need not be re-litigated at every
table. First, the only source of trustworthy graded (per-action-unit, $0/1/2$)
labels is a vet anchor that is \emph{produced}, not downloaded
(Section~\ref{sec:vet-anchor}); this anchor powers only the subordinate, guarded
reliability check and the cited welfare-wrapper plumbing, \emph{not} the
chapter's headline contribution, which is the confound-attribution protocol of
Chapter~\ref{ch:vlm-wrapper} and is validated on planted synthetic controls
without any vet labels. Second, this is a methodology document and we have run
no experiments, so the chapter reports the data \emph{contract} rather than
results. Crucially, the corpus has not yet been forked: the project's
\texttt{data/} substrate is empty at the time of writing, so every corpus count
below --- the $\approx 2040$ records, the $\approx 1547$ distinct images, the
$493$ duplicate filenames, the $\approx 12.7\%$ prevalence, the $336$ source
clips, and the $84$ \texttt{CAT\_01} clips --- is a \emph{projected pre-fork
target} read off the unforked Roboflow source, not a measured or audited
quantity. The figures fix the data contract and the planning targets; they are
never invented results, but neither are they yet measurements.

\section{The forked Roboflow cat-pain corpus --- the binary spine}
\label{sec:roboflow-spine}

\subsection{Purpose and provenance}

The single in-domain corpus is a public Roboflow detection project
(\texttt{lia-k4jkv/cat-pain-ul7lu}) that the project \emph{plans} to fork and
re-host in the workspace as
\texttt{mingraths-workspace/cat-pain-ul7lu-p3rtl}; that fork has not yet been
executed, so the description here is of the intended binary spine rather than an
in-hand corpus. Once forked it will be the only set that lives in the cat-pain
namespace, carrying the full weight of the binary spine: it is to train the
binary pain detector and supply the tight head crops that feed the frozen image
encoder. The unforked source is projected to comprise roughly $2040$ records
that resolve to approximately $1547$ distinct images after de-duplication
(Section~\ref{subsec:hygiene}); these are pre-fork targets, not measured counts. Its upstream provenance is the Zhang/Sun/Tang
2008 Flickr cat-head family of photographs; this is an important honesty point,
because the negative class is therefore a heterogeneous collection of generic
internet cat photographs rather than a clinically curated control group. The
corpus is released under CC~BY~4.0 according to its project metadata, which
makes a derived release possible \emph{provided} it is kept walled from the
non-commercial CatFLW images discussed in
Section~\ref{subsec:catflw} (the license-co-mingling firewall of
Section~\ref{sec:datasheet}).

\subsection{What labels it carries, and what it does not}

The corpus carries binary \texttt{pain}\,/\,\texttt{no\_pain} bounding boxes and
nothing else: there are \emph{zero} per-action-unit, $0$--$10$, or
Feline Grimace Scale (FGS) labels anywhere in it. This is the structural reason
the project cannot derive a graded FGS instrument from this corpus alone, and
it motivates the validated framework of the published FGS literature
\citep{evangelista2019} being used only as background rather than as a label
source. Equally load-bearing is the composition of the negative class: it is an
\emph{unknown mixture} that may include sedated or post-operative animals, and
it is documented as such throughout. The pain labels themselves are best read as
\emph{unvalidated pseudo-labels on generic photographs} until the downstream
confound audit and weak-label reliability pilot have run; nothing in this
chapter treats them as ground truth.

\subsection{Load-bearing data-hygiene facts}
\label{subsec:hygiene}

Table~\ref{tab:hygiene} records the projected hygiene facts that will constrain
every downstream split, count, and claim once the corpus is forked. Each is read
off the unforked Roboflow source rather than measured against a forked export,
and is reproduced from the repository's pre-fork data-contract notes; because
\texttt{data/} is empty at the time of writing, every figure in the table is a
projected pre-fork target, not a measured or audited count. The prevalence
figure deserves special
care: it is pinned to the metadata basis of $264$ pain boxes in $1819$
dataset-wide \texttt{no\_pain} records ($\approx 12.7\%$), with the
live-annotation discrepancy of $246/1811$ ($\approx 12.0\%$) logged rather than
hidden. Critically, the count $264$ is a pain-\emph{box} metadata count and
\emph{not} a count of distinct individuals: the number of distinct painful cats
is not computable without re-identification, which is not performed (the
denominator we report honestly is the distinct-pain-cat count discussed in
Section~\ref{sec:datasheet}). The lone known camera identifier, \texttt{CAT\_01}
(some $84$ clips), is the only handle on individual identity in the entire
corpus.

\begin{table}[t]
\centering
\caption[Load-bearing data-hygiene facts]{Load-bearing data-hygiene facts for
the Roboflow cat-pain corpus. The corpus has not yet been forked and
\texttt{data/} is empty, so every figure is a \emph{projected pre-fork target}
read off the unforked source, not a measured or audited count; each constrains a
downstream gate once the fork is executed.}
\label{tab:hygiene}
\small
\begin{tabular}{@{}p{0.30\linewidth}p{0.40\linewidth}p{0.22\linewidth}@{}}
\toprule
\textbf{Fact} & \textbf{What it means} & \textbf{Consequence} \\
\midrule
$\approx 2040$ records, $\approx 1547$ distinct &
$493$ exact-duplicate filenames projected in the source export &
De-duplicate before splitting (no duplicate may straddle the split) \\
\addlinespace
Prevalence $\approx 12.7\%$ &
$264/1819$ metadata basis; $246/1811$ ($\approx 12.0\%$) live-annotation
discrepancy logged; $1819$ is dataset-wide, not train-split &
$\approx 6.9{:}1$ imbalance; class-weighted loss and oversampling \\
\addlinespace
Leaky shipped split &
$191$ of $336$ source clips ($56.8\%$) appear in both train and valid &
Discard the shipped split entirely; re-split from scratch \\
\addlinespace
$135$ of $336$ clips class-mixed &
A clip cannot be collapsed to a single label &
Stratify at the \emph{image} level, with grouping \\
\addlinespace
\texttt{CAT\_01} lone camera id ($84$ clips) &
The only individual-identity handle in the corpus &
Per-CAT merge validated against \texttt{CAT\_} ids, not pHash/CLIP \\
\addlinespace
Stretch-to-$640$ geometry &
v1 export horizontally inflates $\sim 4{:}3$ faces, distorting AU geometry &
Re-export with Fit (reflect/white edges), augmentation empty \\
\bottomrule
\end{tabular}
\end{table}

The geometry caveat is consequential for a facial-action instrument: the
orbital, muzzle, and whisker action units are geometry-sensitive, and horizontal
stretching of $\sim\!4{:}3$ faces corrupts exactly the signal the ordinal heads
must read. The corpus is therefore re-exported as a fresh immutable version
using a \emph{Fit} (reflect/white edges) resize with no stored augmentation;
augmentation is applied only inside the training loop so that augmented copies
can never enter a reported count. Group identity for cross-validation is parsed
from the filename stem rather than from any parent folder, which neutralises a
phantom collision with an archival cat-head folder; the single numeric
filename collision (clip \texttt{00000100} resolving to two distinct groups) is
asserted explicitly during the parse.

\section{The vet anchor --- the only graded-label source}
\label{sec:vet-anchor}

\subsection{Role and why it is produced, not downloaded}

It is important to be precise about what the vet anchor does \emph{not} do
first. The chapter's headline contribution --- the portable, power-aware
confound-attribution protocol of Chapter~\ref{ch:vlm-wrapper} --- does
\emph{not} depend on the vet anchor at all: it is exercised today on planted
positive and negative synthetic controls. A deletion-safe standalone harness
(\texttt{src/protocols/standalone\_test\_corpus.py}) drives the
no-confound-at-this-power branch, and a planted-positive recovery test
(\texttt{tests/test\_confound\_recovery.py}) asserts that a known
above-threshold capture-condition shift is flagged \texttt{detected}; together
they are where the protocol's behaviour is validated. The vet anchor powers only
two subordinate quantities:
the \emph{demoted, guarded} VLM-as-AU-rater reliability check and the cited
welfare-wrapper plumbing. It is not on the critical path to the headline.

With that ordering fixed, the vet anchor is the single source of trustworthy
per-action-unit $0/1/2$ labels, and it is \emph{produced} rather than downloaded
because no open graded cat-FGS dataset exists. It powers three downstream
quantities and nothing else: the guarded per-AU quadratic weighted kappa of the
VLM-as-AU-rater reliability check (measured against this anchor, and reported
inspected-not-validated with its result pending this independent anchor); the
sensitivity and specificity of the binary decision at the literature-anchored
$0.39$ operating point of the validated FGS \citep{evangelista2019}, estimated
on \emph{these labels only}; and the defer-to-vet abstention curve's one-sided
net-predictive-value lower bound.

These three quantities are not all on equal statistical footing at the budgeted
sample size, and the committed Gate-0 power pre-registration records exactly
which are reportable. The kappa confidence-interval half-width is a budget leg
that the committed sample can meet (the Gate-0 target half-width is $0.15$),
and the abstention net-predictive-value \emph{band} is sized as a one-sided
$95\%$ lower bound and never reported as ``guaranteed''. The $0.39$
point-decision confidence interval, by contrast, is \emph{not} reportable at the
budgeted sample: at the pre-registered $n_{\text{pos}}{=}16$ pain-positive faces
the Clopper--Pearson half-width is $\approx 0.25$, and Gate-0 flags this leg
\texttt{reportable:false}. The $0.39$ sensitivity/specificity point is therefore
carried as \emph{exploratory only}, named as a limitation, and never as a
validation (Chapter~\ref{ch:eval-protocol}). The abstention band itself must be read
with equal care: its current sizing is \emph{placeholder-illustrative}, not a
certified result. The inputs are an unconfirmed placeholder vet budget of $120$
over zero forked data, so the band is reported strictly as an illustrative
one-sided lower bound and is \emph{not} yet certified --- because the budget is a
placeholder over an empty corpus, not because the arithmetic forbids a useful
floor. The size of the anchor is itself not a guess: it is the \emph{single
integer} fixed by the upstream power-calculation gate before any veterinary hour
is spent. The committed pre-gate figure is on the order of $120$ confirmed
images with at least $50$ pain-positive, matching the Gate-0 power
pre-registration; the larger ``$150$'' often quoted informally is not
power-backed and is not used here. The committed integer replaces this
placeholder once a real budget is set over real data, and only a non-placeholder
budget over measured faces can move the abstention band from illustrative to
certified.

\subsection{What it carries and how it is produced}

The anchor carries five action units --- ear, orbital, muzzle, whiskers, and
head --- each scored $0/1/2$ (absent / moderate-or-uncertain / obvious, in the
validated-instrument wording of \citealp{evangelista2019}) and confirmed by a
veterinarian. The $0$--$10$ sum and the $0.39$ decision ratio are computed in
code and are \emph{never} elicited from the labeler. The anchor is produced by a
review loop: a vision-language model pre-fills a forced-enumeration structured
output (rationale before score, one $0/1/2$ field per AU); cases are triaged to
the vet by value-per-hour (pain-positive, near-threshold, the weakest AUs, and
confident-learning flags); the vet accepts or corrects prefilled scores in a
review interface rather than labelling from scratch; and the per-AU kappa is then
computed row-paired with quadratic weights. Because correction is faster than
from-scratch labelling, the anchor is the project's binding cost in vet hours
rather than in money or compute.

\subsection{The circularity firewall}
\label{subsec:firewall}

A non-negotiable firewall governs how the anchor is used. Sensitivity and
specificity at the $0.39$ operating point are estimated \emph{only} on these
vet-confirmed labels. The kappa of the model's predictions against the VLM is
\emph{never} reported as validation, because it would measure the model
re-learning the VLM's heuristic rather than agreement with clinical ground
truth. The VLM-versus-vet kappa is a \emph{method whose result is pending} the
independent per-AU vet anchor --- it ships today as a runnable protocol for
asking whether a frozen VLM can weak-label feline FGS action units at human-rater
agreement, not as a number, since the binary corpus cannot yield the $0/1/2$
ground truth the kappa needs --- and is in no case a validation of the decision.
A high kappa would measure weak-labeling \emph{capability} only if the anchor is
independent of the rater who scored the target \emph{and} the rubric handed to
the VLM differs from the one the vet scored from; otherwise it measures
rubric-following, and a run must state which it measures
(see \texttt{src/eval/kappa.py}). The kappa check is, throughout, a
\emph{guarded subordinate reliability check} and not the contribution: its
confidence-interval lower-bound acceptance gate is standard clinimetrics rather
than anything novel here, and it serves as kill-tree insurance for the headline
protocol, not as the headline. The graded $0$--$10$ layer ships
\emph{inspected-not-validated}: the word ``graded'' is struck from every
validated-claim sentence, and the headline confound-attribution protocol holds
with both the kappa check and the graded layer dropped.

\section{Scaffolding-only sources}
\label{sec:scaffolding}

The remaining sources are scaffolding, preprocessing, or unit-test fixtures.
None of them produces a headline number; each is walled in its own namespace so
that a non-cat or non-commercial image can never reach a cat-pain loader. Their
inventory and namespacing are summarised in Table~\ref{tab:inventory} at the end
of the chapter.

\subsection{Horse-grimace --- decode unit-test fixture only}
\label{subsec:horse}

The horse-grimace region set
(\texttt{oliveirabruno01/openfarm-horse-grimace-region}, CC~BY~4.0) is the only
open AU-graded grimace set in the required $0/1/2$ shape, and it is used
\emph{exclusively} to unit-validate the decode path
(CORN-decode $\rightarrow$ per-AU $0$--$2$ $\rightarrow$ $0$--$10$ sum
$\rightarrow$ $0.39$ threshold) on real graded labels, and optionally to
warm-start three of the five heads. Its caveats are baked in: only five unique
horses and three of five AUs (no whiskers, no head), which makes it a coarse
ordinal-direction prior and \emph{never} calibrated transfer. No horse number
appears in the abstract, results, or any transfer table; the only permissible
sentence is the methods-appendix note that the decode path was unit-validated on
five-horse genuine $0/1/2$ labels.

\subsection{CatFLW --- eye-alignment landmarks only}
\label{subsec:catflw}

The CatFLW landmark set (\texttt{georgemartvel/catflw}; some $2079$ cat faces
with $48$ CatFACS-aligned landmarks and a bounding box, but no pain or AU
scores) is used solely for the alignment / normalised-mean-error audit that
decides whether a detector box crop is geometrically clean enough to feed the
ordinal heads, and to supply a two-point eye-similarity transform when it is not
\citep{martvel2024catflw}. Using it as a graded anchor would be a category
error, and it is explicitly cut from the critical path. Its license is the
binding constraint: CatFLW is CC~BY-NC~4.0, so any aligner trained on it inherits
the non-commercial restriction, and these images \emph{must not} leak into the
CC~BY release built from the spine corpus.

\subsection{Frozen encoder, fallback cropper, and archival aligner}

Three further assets are conceded plumbing. The frozen image encoder
(\texttt{facebook/dinov2-small}, Apache-2.0) is a foundation-feature backbone
\citep{oquab2023dinov2} that is never domain-adapted; its features are cached to
disk once and the ordinal heads are trained off the cache. The register variant
(\texttt{dinov2\_vits14\_reg}) is the field default --- registers suppress
attention artifacts that hurt the dense, localized features the per-AU heads read
(orbital, ear, muzzle sub-regions) --- with the plain backbone kept only as an
A/B comparator before \mbox{ViT-S/14} is locked; this remains conceded plumbing,
not a contribution. A Haar
\texttt{frontalcatface} cascade (MIT) is a zero-training fallback cropper used
only when a detector box is missing, and a detection failure is logged as an
abstention signal rather than promoted into a landmark pipeline. The Zhang-2008
archive (CC0) is namespaced \emph{outside} the cat-pain directory and serves only
as a fallback two-point eye-aligner if CatFLW access stalls. None of these
assets is claimed as a contribution, consistent with the project's standing
concession that the encoder-plus-ordinal-head engine is plumbing. This framing
also keeps the project distinct from prior landmark-and-classifier feline-pain
work that this corpus could otherwise be made to replicate
\citep{feighelstein2023}.

\section{Data-provenance datasheet and the license firewall}
\label{sec:datasheet}

The released artifact carries a provenance datasheet that is the project's
honesty layer; its co-mingling firewall and namespacing flow are shown in
Figure~\ref{fig:provenance}, and the full source inventory in
Table~\ref{tab:inventory}. The firewall has two jobs. First, it keeps each
license in its own namespace, so that the CC~BY-NC CatFLW images and any aligner
derived from them never reach a CC~BY release built from the spine corpus, and
so that horse, mouse, or archival cat-head images can never enter a cat-pain
loader. Second, it enforces the per-CAT provenance discipline: clips are merged
toward true individuals \emph{only} to assert fold disjointness and to print the
honest distinct-pain-cat denominator, never to fuse the lone \texttt{CAT\_01}
camera into a single degenerate cross-validation group. Anti-benchmark
discipline is part of the same layer: augmented copies never enter a reported
count, every reported proportion is accompanied by a distinct-pain-cat
denominator and Clopper--Pearson or bootstrap confidence intervals, and the
framing is decision-support triage --- never an autonomous analgesia trigger.

\begin{figure}[t]
\centering
\begin{tikzpicture}[
  node distance=8mm and 12mm,
  box/.style={draw, rounded corners, align=center, font=\small,
              minimum height=8mm, inner sep=4pt, fill=black!3},
  spine/.style={box, fill=blue!6},
  anchorbox/.style={box, fill=green!8},
  scaffold/.style={box, fill=orange!7},
  flow/.style={-{Latex}, thick},
  note/.style={font=\scriptsize\itshape, align=center}
]
% spine pipeline (top row)
\node[spine] (fork)   {Forked Roboflow\\cat-pain (CC BY 4.0)\\$\approx 2040$ records};
\node[spine, right=of fork] (dedup) {De-dup\\($493$ dup names)\\$\rightarrow \approx 1547$};
\node[spine, right=of dedup] (merge) {Per-CAT merge\\(vs \texttt{CAT\_} ids)\\disjointness only};
\node[spine, right=of merge] (holdout) {Frozen, hashed\\cat-disjoint\\hold-out};
% vet anchor branch (bottom)
\node[anchorbox, below=14mm of merge] (vlm) {VLM pre-fill\\(5 AUs, enum $0/1/2$)};
\node[anchorbox, left=of vlm] (vet) {Vet accept/correct\\$=$ \textbf{vet anchor}\\(only graded source)};
\node[anchorbox, right=of vlm] (kappa) {$\kappa$ vs vet\\+ $0.39$ sens/spec\\(firewalled)};
% scaffolding (far left, walled off)
\node[scaffold, below=10mm of fork] (scaff) {Scaffolding namespaces\\
  horse (CC BY) $\cdot$ CatFLW (CC BY-NC)\\
  DINOv2 (Apache) $\cdot$ Haar (MIT) $\cdot$ Zhang (CC0)};
% spine flow
\draw[flow] (fork) -- (dedup);
\draw[flow] (dedup) -- (merge);
\draw[flow] (merge) -- (holdout);
% anchor flow
\draw[flow] (vet) -- (vlm);
\draw[flow] (vlm) -- (kappa);
\draw[flow] (merge) -- node[note, right] {crops} (vlm);
% firewall: scaffolding cannot reach spine
\draw[red, very thick, dashed] ($(scaff.north west)+(-0.2,0.35)$)
   -- ($(scaff.north east)+(2.6,0.35)$)
   node[midway, above, note, text=red] {license / namespace firewall (no co-mingling)};
\end{tikzpicture}
\caption[Data-provenance and namespacing firewall]{Data-provenance and
namespacing firewall. The spine corpus (blue) is de-duplicated, merged per CAT
\emph{only} to assert disjointness, and frozen into a hashed cat-disjoint
hold-out; the vet-anchor branch (green) is the only graded-label source and its
kappa-versus-vet and sensitivity/specificity outputs are firewalled per
Section~\ref{subsec:firewall}; scaffolding sources (orange) are walled in their
own namespaces and never reach a cat-pain loader or a CC~BY release.}
\label{fig:provenance}
\end{figure}

\begin{table}[t]
\centering
\caption[Source inventory and namespacing]{Source inventory: directory, role,
license, and whether the source carries AU labels. Only the vet anchor carries
graded ($0/1/2$) AU labels; only the spine and vet anchor are in the cat-pain
namespace.}
\label{tab:inventory}
\small
\begin{tabular}{@{}p{0.24\linewidth}p{0.21\linewidth}p{0.27\linewidth}p{0.20\linewidth}@{}}
\toprule
\textbf{Source} & \textbf{Role} & \textbf{License} & \textbf{Carries AU labels?} \\
\midrule
Forked Roboflow cat-pain & Binary spine (detector + crop source) & CC~BY~4.0 &
No (binary boxes only) \\
\addlinespace
Vet anchor (size fixed by power gate) & Only graded source (guarded $\kappa$ check; $0.39$ sens/spec estimate, non-reportable at $n_{\text{pos}}{=}16$) &
Project-original (derived) & \textbf{Yes} (per-AU $0/1/2$) \\
\addlinespace
Horse-grimace & Decode unit-test fixture & CC~BY~4.0 &
$3/5$ AUs, $5$ horses (fixture only) \\
\addlinespace
CatFLW & Eye-alignment / NME audit & CC~BY-NC~4.0 &
No ($48$ landmarks, no pain) \\
\addlinespace
\texttt{dinov2-small} & Frozen encoder (plumbing) & Apache-2.0 & n/a \\
\addlinespace
Haar \texttt{frontalcatface} & Fallback cropper / abstention signal & MIT & n/a \\
\addlinespace
Zhang-2008 archive & Fallback aligner (if CatFLW blocked) & CC0 &
No ($9$ coarse landmarks) \\
\bottomrule
\end{tabular}
\end{table}

Finally, Figure~\ref{fig:crops} is a captioned placeholder for the qualitative
crop-acceptance illustration to be produced once the alignment audit has run; no
example images are fabricated here, consistent with the honesty rule that
governs this document.

\begin{figure}[t]
\centering
\fbox{\parbox[c][0.16\textheight][c]{0.86\linewidth}{\centering
\emph{Placeholder --- figure to be produced.}\\[2pt]
Example acceptable (frontal, eye-aligned, sufficient face-pixel resolution)
versus excluded (profile, eyes-closed, occluded, or below the
normalised-mean-error threshold) crops, drawn from the alignment / NME audit of
Section~\ref{subsec:catflw}. Excluded crops are routed to the vet and logged as
abstention signals; no crop images are fabricated in this methodology document.}}
\caption[Acceptable versus excluded crops (placeholder)]{Acceptable versus
excluded face crops (to be produced). The acceptance criterion is the geometric
quality gate of the alignment audit, not pain content.}
\label{fig:crops}
\end{figure}
